\documentclass{article}

\usepackage{amsmath,amssymb}
\usepackage{xurl}
\usepackage[colorlinks=true,linkcolor=blue,citecolor=blue,urlcolor=blue]{hyperref}

\input{./command}

\title{Global and Per-Sector Unit-Weight Witnesses in the CPSV16 BNT Argument}
\author{}
\date{2026-05-30}

\begin{document}
\maketitle
\gapnote{scope-restriction}{resolved}

\section{Closure record}

Cirac--P\'erez-Garc\'ia--Schuch--Verstraete normalize the raw coefficients in
the basis-of-normal-tensors expansion by one global condition. Section~II.C,
line~246 of Ref.~\cite{Cirac2016MPDO_arXiv} requires
$|\mu_{j,q}|\leq 1$ for every copy weight and
$|\mu_{j_\ast,q_\ast}|=1$ for at least one copy. The paper-level canonical
form does not require every BNT sector to contain a unit-modulus copy. One
sector may contain only strictly subunit copy weights while another sector
contains the global unit witness. The formal predicate
\leanid{MPSTensor.IsBNTCanonicalForm} records exactly this global condition
through the field
\leanid{MPSTensor.IsBNTCanonicalForm.weight_unit_exists}.

Earlier full-basis matching and global-gauge theorems assumed the stronger
per-sector condition $\forall j,\ \exists q,\ |\mu_{j,q}|=1$. That condition
entered through an asymptotic matching argument that projected onto individual
sectors using the Ces\`aro non-decay lemma. The restriction has been
eliminated.

For sector $j$, define the length-$N$ coefficient sum by
\begin{align}
    c_N^{(j)}
        &= \sum_q \mu_{j,q}^N.
    \label{eq:cpsvunit_sector_coefficient_sum}
\end{align}
The exact linear-independence matcher works at one fixed sufficiently large
length. Its equal-case declaration is
\leanid{MPSTensor.exists_block_match_exact}, and its proportional-case
declaration is
\leanid{MPSTensor.exists_block_match_exact_of_eventuallyProportional}. The
theorem
\leanid{MPSTensor.IsBNTCanonicalForm.coeff_not_eventually_zero} shows that
$c_N^{(j)}$ is nonzero for infinitely many $N$ whenever the copy weights are
nonzero; it requires no modulus assumption. The matching and global-gauge
declarations therefore use exactly the source hypothesis set.

The unitary equal-case assembly in Corollary~A.6 of
Ref.~\cite{Cirac2016MPDO_arXiv}, identified in the source by
\texttt{thm:Fundamental-CFII} at
\path{Papers/1606.00608/MPDO-22-12-17-2.tex:1197}, is complete under the same
hypotheses. Its formal declarations are
\leanid{MPSTensor.ft_sector_bnt_equal_unitary_global_gauge_witnesses} and
\leanid{MPSTensor.fundamentalTheorem_equal_canonicalForm_unitary}. Neither
declaration assumes a per-sector unit weight.

The faithful structural normalization remains the global unit-weight field
\path{MPSTensor.IsBNTCanonicalForm.weight_unit_exists} in
\path{TNLean/MPS/FundamentalTheorem/SectorBNT/Basic.lean}. The interaction
between this convention and the scale fixed by the MPDO renormalization
equations is treated separately in
\path{docs/paper-gaps/cpsv16_unit_weight_rfp_scale_tension.tex}.

\bibliographystyle{alpha}
\bibliography{references}

\end{document}
